% GENERATED FILE. Edit canonical lessons, then run npm run generate:books.
% canonical-source-hash: fnv1a64:9199af5ef0f740e2
% canonical-lessons: UR-C430-para, UR-C430-digri, UR-C430-tapna, UR-C430-habs, UR-C430-jhalna

\chapter{Heat}
\label{ch:ur-heat}
\hlchaptermodality{430}

\begin{chapteropening}
\textbf{By the end of this chapter:} \emph{I can talk about a hot day: the temperature, the airless heat and keeping cool.}

Complete the last lesson of chapter 430: I can talk about a hot day: the temperature, the airless heat and keeping cool.
\end{chapteropening}

\section[\texorpdfstring{\ur{پارہ} (pārā)}{پارہ (pārā)}]{\ur{پارہ} (pārā) — the temperature (in everyday talk); mercury}

\label{lesson:UR-C430-para}



\begin{quote}
Before the new one: say the Urdu for to shrink, then the Urdu for to go stiff.
\end{quote}

\subsection*{You'll want to know: \ur{پارہ}}
\textbf{\ur{پارہ}} — \emph{pārā} — \textquotedblleft{}the temperature (in everyday talk); mercury\textquotedblright{}.

\begin{quote}
\ur{گرمیوں} \ur{میں} \ur{پارہ} \ur{چالیس} \ur{سے} \ur{اوپر} \ur{جاتا} \ur{ہے۔}
\end{quote}

\begin{quote}
\emph{garmiyoṅ meṅ pārā chālīs se ūpar jātā hai.}
\end{quote}

\begin{quote}
In summer the temperature goes above forty.
\end{quote}

Masculine.

\begin{cousinweb}
\textbf{\ur{پارہ}} (\emph{pārā}) comes from Sanskrit \emph{pārada}, \textquotedblleft{}quicksilver\textquotedblright{} (Monier-Williams); the mercury in a thermometer gave Urdu its everyday word for the temperature.
\end{cousinweb}

\begin{grammarlens}[title={What you've built}]
One more word for getting around.
\end{grammarlens}

\subsection*{Your turn}
\begin{itemize}
\raggedright
  \item \emph{Say it:} \emph{pārā}
  \item \emph{Say it:} \emph{pārā}, once more
  \item \emph{Say it:} \emph{garmiyoṅ meṅ pārā chālīs se ūpar jātā hai.} — In summer the temperature goes above forty.
  \item \emph{Recall:} say the Urdu for to shrink, then the Urdu for to go stiff, then say \emph{pārā} again
\end{itemize}

\begin{tcolorbox}[breakable,colback=teal!4,colframe=teal!35!black,title={Before you move on}]
What does \ur{پارہ} mean? (\textquotedblleft{}The temperature (in everyday talk); mercury\textquotedblright{}.) Say it once more.
\end{tcolorbox}
\section[\texorpdfstring{\ur{ڈگری} (ḍigrī)}{ڈگری (ḍigrī)}]{\ur{ڈگری} (ḍigrī) — a degree (of temperature)}

\label{lesson:UR-C430-digri}



\begin{quote}
Before the new one: say the Urdu for to go stiff, then the Urdu for the temperature (in everyday talk).
\end{quote}

\subsection*{You'll want to know: \ur{ڈگری}}
\textbf{\ur{ڈگری}} — \emph{ḍigrī} — \textquotedblleft{}a degree (of temperature)\textquotedblright{}.

\begin{quote}
\ur{آج} \ur{چالیس} \ur{ڈگری} \ur{گرمی} \ur{ہے۔}
\end{quote}

\begin{quote}
\emph{āj chālīs ḍigrī garmī hai.}
\end{quote}

\begin{quote}
It's forty degrees today.
\end{quote}

Feminine.

\begin{cousinweb}
\textbf{\ur{ڈگری}} (\emph{ḍigrī}) is borrowed from English \emph{degree}.
\end{cousinweb}

\begin{grammarlens}[title={What you've built}]
One more word for getting around.
\end{grammarlens}

\subsection*{Your turn}
\begin{itemize}
\raggedright
  \item \emph{Say it:} \emph{ḍigrī}
  \item \emph{Say it:} \emph{ḍigrī}, once more
  \item \emph{Say it:} \emph{āj chālīs ḍigrī garmī hai.} — It's forty degrees today.
  \item \emph{Recall:} say the Urdu for to go stiff, then the Urdu for the temperature (in everyday talk), then say \emph{ḍigrī} again
\end{itemize}

\begin{tcolorbox}[breakable,colback=teal!4,colframe=teal!35!black,title={Before you move on}]
What does \ur{ڈگری} mean? (\textquotedblleft{}A degree (of temperature)\textquotedblright{}.) Say it once more.
\end{tcolorbox}
\section[\texorpdfstring{\ur{تپنا} (tapnā)}{تپنا (tapnā)}]{\ur{تپنا} (tapnā) — to be burning hot; to scorch}

\label{lesson:UR-C430-tapna}



\begin{quote}
Before the new one: say the Urdu for the temperature (in everyday talk), then the Urdu for a degree (of temperature).
\end{quote}

\subsection*{You'll want to know: \ur{تپنا}}
\textbf{\ur{تپنا}} — \emph{tapnā} — \textquotedblleft{}to be burning hot; to scorch\textquotedblright{}.

\begin{quote}
\ur{دوپہر} \ur{کو} \ur{سڑک} \ur{دھوپ} \ur{میں} \ur{تپتی} \ur{ہے۔}
\end{quote}

\begin{quote}
\emph{dopahar ko saṛak dhūp meṅ taptī hai.}
\end{quote}

\begin{quote}
At midday the road bakes in the sun.
\end{quote}

Past (it …): \ur{وہ} \ur{زمین} \ur{دن} \ur{بھر} \ur{تپی}, \emph{voh zamīn din bhar tapī}.

\begin{cousinweb}
\textbf{\ur{تپنا}} (\emph{tapnā}) is inherited from Sanskrit \emph{tápyati}, \textquotedblleft{}is hot\textquotedblright{}: CDIAL lists the Hindi cousin \emph{tāpnā}, \textquotedblleft{}to warm oneself\textquotedblright{}, under it (CDIAL 5684).
\end{cousinweb}

\begin{grammarlens}[title={What you've built}]
One more word for getting around.
\end{grammarlens}

\subsection*{Your turn}
\begin{itemize}
\raggedright
  \item \emph{Say it:} \emph{tapnā}
  \item \emph{Say it:} \emph{tapnā}, once more
  \item \emph{Say it:} \emph{dopahar ko saṛak dhūp meṅ taptī hai.} — At midday the road bakes in the sun.
  \item \emph{Recall:} say the Urdu for the temperature (in everyday talk), then the Urdu for a degree (of temperature), then say \emph{tapnā} again
\end{itemize}

\begin{tcolorbox}[breakable,colback=teal!4,colframe=teal!35!black,title={Before you move on}]
What does \ur{تپنا} mean? (\textquotedblleft{}To be burning hot; to scorch\textquotedblright{}.) Say it once more.
\end{tcolorbox}
\section[\texorpdfstring{\ur{حبس} (habs)}{حبس (habs)}]{\ur{حبس} (habs) — airless heat; sultriness}

\label{lesson:UR-C430-habs}



\begin{quote}
Before the new one: say the Urdu for a degree (of temperature), then the Urdu for to be burning hot.
\end{quote}

\subsection*{You'll want to know: \ur{حبس}}
\textbf{\ur{حبس}} — \emph{habs} — \textquotedblleft{}airless heat; sultriness\textquotedblright{}.

\begin{quote}
\ur{بارش} \ur{سے} \ur{پہلے} \ur{بڑا} \ur{حبس} \ur{تھا۔}
\end{quote}

\begin{quote}
\emph{bārish se pehle baṛā habs thā.}
\end{quote}

\begin{quote}
Before the rain it was very close and airless.
\end{quote}

Masculine.

\begin{cousinweb}
\textbf{\ur{حبس}} (\emph{habs}) is usually traced, through Persian, to Arabic \emph{ḥabs}, \textquotedblleft{}shutting in, confinement\textquotedblright{}; the sense \textquotedblleft{}airless heat\textquotedblright{} probably grew in Urdu.
\end{cousinweb}

\begin{grammarlens}[title={What you've built}]
One more word for getting around.
\end{grammarlens}

\subsection*{Your turn}
\begin{itemize}
\raggedright
  \item \emph{Say it:} \emph{habs}
  \item \emph{Say it:} \emph{habs}, once more
  \item \emph{Say it:} \emph{bārish se pehle baṛā habs thā.} — Before the rain it was very close and airless.
  \item \emph{Recall:} say the Urdu for a degree (of temperature), then the Urdu for to be burning hot, then say \emph{habs} again
\end{itemize}

\begin{tcolorbox}[breakable,colback=teal!4,colframe=teal!35!black,title={Before you move on}]
What does \ur{حبس} mean? (\textquotedblleft{}Airless heat; sultriness\textquotedblright{}.) Say it once more.
\end{tcolorbox}
\section[\texorpdfstring{\ur{جھلنا} (jhalnā)}{جھلنا (jhalnā)}]{\ur{جھلنا} (jhalnā) — to fan (someone)}

\label{lesson:UR-C430-jhalna}



\begin{quote}
Before the new one: say the Urdu for to be burning hot, then the Urdu for airless heat.
\end{quote}

\subsection*{You'll want to know: \ur{جھلنا}}
\textbf{\ur{جھلنا}} — \emph{jhalnā} — \textquotedblleft{}to fan (someone)\textquotedblright{}.

\begin{quote}
\ur{گرمی} \ur{میں} \ur{ماں} \ur{نے} \ur{بچے} \ur{کو} \ur{پنکھا} \ur{جھلا۔}
\end{quote}

\begin{quote}
\emph{garmī meṅ māṅ ne bachche ko pankhā jhalā.}
\end{quote}

\begin{quote}
In the heat the mother fanned the child.
\end{quote}

Past, with \ur{نے} (he or she …): \ur{اس} \ur{نے} \ur{پنکھا} \ur{جھلا}, \emph{us ne pankhā jhalā}.

\begin{cousinweb}
\textbf{\ur{جھلنا}} (\emph{jhalnā}) is inherited: CDIAL lists the Hindi cousin \emph{jhalnā}, \textquotedblleft{}to move a fan to and fro\textquotedblright{}, under \emph{jhal}, \textquotedblleft{}a sudden movement\textquotedblright{}, a form it reconstructs (CDIAL 5351).
\end{cousinweb}

\begin{grammarlens}[title={What you've built}]
One more word for getting around.
\end{grammarlens}

\subsection*{Your turn}
\begin{itemize}
\raggedright
  \item \emph{Say it:} \emph{jhalnā}
  \item \emph{Say it:} \emph{jhalnā}, once more
  \item \emph{Say it:} \emph{garmī meṅ māṅ ne bachche ko pankhā jhalā.} — In the heat the mother fanned the child.
  \item \emph{Recall:} say the Urdu for to be burning hot, then the Urdu for airless heat, then say \emph{jhalnā} again
\end{itemize}

\begin{tcolorbox}[breakable,colback=teal!4,colframe=teal!35!black,title={Before you move on}]
What does \ur{جھلنا} mean? (\textquotedblleft{}To fan (someone)\textquotedblright{}.) Say it once more.
\end{tcolorbox}
